% TEMPLATE-MILC-v3.tex - plantilla maestra UNICA de las guias MILC v3.
%
% La usan las tres familias de guias, cada una con su driver:
%   · Grados 6-11  -> scripts/build-guias-g11.py
%   · Bebras       -> scripts/build-guias-bebras.py
%   · Semillero    -> scripts/build-guias-semillero.py
%
% Antes habia tres copias casi identicas (~400 lineas iguales, 6 distintas):
% cada mejora habia que aplicarla tres veces, y en la practica no se hacia
% (el motor de recursos vivia solo en dos de ellas). Lo que varia entre
% familias esta parametrizado en 6 placeholders que cada driver llena
% (se nombran aqui SIN su sintaxis real: el builder reemplaza por texto plano,
% tambien dentro de los comentarios, y un valor multilinea rompe el preambulo):
%
%   HEADER_IZQ        encabezado izquierdo de pagina
%   HEADER_DER        encabezado derecho de pagina
%   PORTADA_ETIQUETA  rotulo sobre el numero grande (GRADO/MOMENTO/MODULO)
%   PORTADA_SERIE     linea de serie de la portada
%   PORTADA_ID        identificador de la guia en la portada
%   PORTADA_CIRCULO   el circulito de grado (°), o vacio si no aplica
%
% El resto de claves las documenta content/guias/_SCHEMA.md. El script valida
% que no queden placeholders sin reemplazar antes de escribir el archivo final.

\documentclass[11pt,letterpaper]{article}
\usepackage[margin=2.0cm]{geometry}
\usepackage[table]{xcolor}
\usepackage{fontspec}
\usepackage[spanish]{babel}
\usepackage{tikz}
% Librerías TikZ para los diagramas de las guías (bloque `recursos:`):
%  · babel  → babel en español vuelve activos caracteres como «>», y TikZ los usa
%             en sus opciones (>=stealth, ->). Sin esta librería, cualquier
%             diagrama con flechas revienta con «\language@active@arg> has an
%             extra }». Debe ir ANTES de que se dibuje nada.
%  · positioning        → sintaxis «below left=2pt and 3pt».
%  · decorations.pathreplacing → llaves ({brace}) para señalar rangos.
\usetikzlibrary{babel,positioning,decorations.pathreplacing}
\usepackage{graphicx}
% Circuitos: el trabajo de electrónica se hace hoy con micro:bit y sus sensores
% integrados (MakeCode, sin cableado), así que no se necesitan esquemáticos.
% Si algún día entran montajes con protoboard, basta descomentar esta línea:
% los diagramas .tex/.tikz declarados en `recursos:` ya se incrustan solos.
% \usepackage{circuitikz}
\usepackage[most]{tcolorbox}
\usepackage{tabularx}
\usepackage{array}
\usepackage{fancyhdr}
\usepackage{titlesec}
\usepackage[document]{ragged2e}  % bandera a la derecha en todo el cuerpo
\usepackage{enumitem}
\usepackage{microtype}

% Helvetica Neue no trae la flecha U+2192, asi que cada «→» escrito literalmente
% en el YAML se compilaba a NADA, en silencio (462 flechas en 61 guias: rutas del
% tipo «paleta → tipografia → lockup» salian rotas). Mapeamos el caracter a su
% version matematica, que si tiene glifo. Asi el autor puede seguir escribiendo
% «→» comodamente en el YAML.
\usepackage{newunicodechar}
\newunicodechar{→}{\ensuremath{\rightarrow}}
% Mismo problema con los simbolos de marca y vinetas: el «✓» de «una palomita ✓»
% salia como cuadrito vacio en el PDF de todas las guias que lo usaban.
\usepackage{pifont}
\newunicodechar{✓}{\ding{51}}
\newunicodechar{✗}{\ding{55}}
\newunicodechar{✔}{\ding{52}}
\newunicodechar{•}{\textbullet}
\newunicodechar{≥}{\ensuremath{\geq}}
\newunicodechar{≤}{\ensuremath{\leq}}

\setmainfont{Helvetica Neue}
\setsansfont{Helvetica Neue}
\setlength{\parindent}{0pt}
\setlength{\parskip}{6pt}
% Sin particion de palabras: en 6.º-7.º una palabra cortada por guion al final
% de linea («Comuni-cacion») frena la lectura mucho mas que una linea corta.
\hyphenpenalty=10000
\exhyphenpenalty=10000
\pretolerance=2000
\tolerance=3000
\setlength{\emergencystretch}{3em}
\linespread{1.10}
\renewcommand{\arraystretch}{1.25}
\setlist{leftmargin=5mm,itemsep=1mm,topsep=1mm}
\newcolumntype{Y}{>{\RaggedRight\arraybackslash}X}

\definecolor{milcMagenta}{HTML}{E6007E}
\definecolor{milcVino}{HTML}{5A0038}
\definecolor{milcTurquesa}{HTML}{0093A5}
\definecolor{milcVerde}{HTML}{58A923}
\definecolor{milcMostaza}{HTML}{E5B400}
\definecolor{milcOcre}{HTML}{B86600}
\definecolor{milcNegro}{HTML}{241816}
\definecolor{milcGris}{HTML}{F7F7F4}
\definecolor{milcLinea}{HTML}{E8E2DD}
\definecolor{milcGradePrimary}{HTML}{0D9488}
\definecolor{milcGradeDark}{HTML}{115E59}
\definecolor{milcGradeSoft}{HTML}{CCFBF1}
\definecolor{milcGradeAccent}{HTML}{CFE7FF}
\definecolor{milcGradeText}{HTML}{FFFFFF}
\color{milcNegro}

\pagestyle{fancy}
\fancyhf{}
\lhead{\small Territorio interior · Educación socioemocional}
\rhead{\small Grado 8° · Momento 9 · MILC v3}
\cfoot{\small\thepage}
\renewcommand{\headrulewidth}{0pt}

\newtcolorbox{titlebox}[1]{
  enhanced,colback=#1,colframe=#1,arc=7mm,boxrule=0pt,
  left=7mm,right=7mm,top=3mm,bottom=3mm,
  before skip=8pt,after skip=8pt,
  fontupper=\sffamily\bfseries\Large\color{white}
}

\newtcolorbox{softbox}[3]{
  enhanced,colback=#2,colframe=#1,borderline west={4pt}{0pt}{#1},
  arc=2mm,boxrule=.4pt,left=4mm,right=4mm,top=3mm,bottom=3mm,
  before skip=6pt,after skip=8pt,title={#3},coltitle=white,
  colbacktitle=#1,fonttitle=\sffamily\bfseries,fontupper=\RaggedRight,
  attach boxed title to top left={xshift=3mm,yshift=-2mm},
  boxed title style={arc=2mm, boxrule=0pt}
}

\newtcolorbox{drawbox}[2]{
  enhanced,colback=white,colframe=#1,arc=4mm,boxrule=1pt,
  left=5mm,right=5mm,top=4mm,bottom=4mm,before skip=8pt,after skip=8pt,
  title={#2},coltitle=white,colbacktitle=#1,fonttitle=\sffamily\bfseries,
  fontupper=\RaggedRight,
  attach boxed title to top left={xshift=4mm,yshift=-2mm},
  boxed title style={arc=3mm, boxrule=0pt}
}

\newtcolorbox{infoband}[1]{
  enhanced,colback=#1!8,colframe=#1!50,arc=2mm,boxrule=.5pt,
  left=4mm,right=4mm,top=2mm,bottom=2mm,before skip=4pt,after skip=4pt,
  fontupper=\small\RaggedRight
}

% Puente narrativo entre secciones (contrato editorial Conéctate).
% Da continuidad: "Acabas de X. Ahora vas a Y porque Z."
% Visualmente: banda izquierda en color del grado, texto en itálica pequeña.
\newtcolorbox{puentebox}{
  enhanced,colback=milcGradeSoft,colframe=milcGradeDark,
  borderline west={3pt}{0pt}{milcGradeDark},
  arc=0mm,boxrule=0pt,left=5mm,right=4mm,top=2.5mm,bottom=2.5mm,
  before skip=4pt,after skip=8pt,
  fontupper=\itshape\small\color{milcNegro}\RaggedRight
}
% La flecha va en modo matematico a proposito: Helvetica Neue NO tiene el glifo
% U+2192, asi que \textrightarrow se compilaba a NADA («Missing character: There
% is no → in font Helvetica Neue») y el puente salia sin flecha. En modo
% matematico la flecha viene de la fuente matematica y si se dibuja.
\newcommand{\puente}[1]{%
  \begin{puentebox}%
    \textcolor{milcGradeDark}{$\rightarrow$}\hspace{2mm}#1%
  \end{puentebox}%
}

\newcommand{\linead}[1]{\noindent\color{milcLinea}\rule{#1}{0.5pt}\color{milcNegro}}
\newcommand{\respuesta}[1]{\par\vspace{2mm}\foreach \n in {1,...,#1}{\noindent\color{milcLinea}\rule{\linewidth}{0.5pt}\par\vspace{5mm}}\color{milcNegro}}
\newcommand{\checkbox}{\textcolor{milcGradeDark}{\fbox{\phantom{x}}}\hspace{5pt}}

% ─── Listas aireadas para las guías socioemocionales ────────────────────
% Componer las verificaciones como «\checkbox texto\\» deja el texto que salta
% de línea debajo del cuadrito y apelmaza el bloque. Estos entornos alinean la
% continuación y airean los ítems. Los usa Territorio Interior; cualquier
% familia puede hacerlo.
\newlist{checklista}{itemize}{1}
\setlist[checklista]{
  label=\textcolor{milcGradeDark}{\fbox{\phantom{x}}},
  leftmargin=9mm, labelsep=3.2mm, itemsep=2.4mm,
  topsep=2.5mm, parsep=0pt, partopsep=0pt, before=\RaggedRight
}
\newlist{pasoslista}{enumerate}{1}
\setlist[pasoslista]{
  label=\textcolor{milcGradeDark}{\sffamily\bfseries\arabic*.},
  leftmargin=8mm, labelsep=3mm, itemsep=2.8mm,
  topsep=1mm, parsep=0pt, partopsep=0pt, before=\RaggedRight
}

\newcommand{\phasecard}[4]{
\begin{tcolorbox}[enhanced,colback=#3,colframe=#2,arc=3mm,boxrule=.5pt,height=3.05cm,valign=center,halign=center,left=1.5mm,right=1.5mm,top=2mm,bottom=2mm]
{\sffamily\bfseries\fontsize{22}{24}\selectfont\textcolor{#2}{#1}}\par
{\sffamily\bfseries\small #4}
\end{tcolorbox}
}

% ── Piezas del rediseno 2026 (legibilidad 6.º-11.º) ───────────────────
% Banda de estacion: reemplaza el titlebox macizo. Titulo a la izquierda,
% dato practico (tiempo/modalidad) a la derecha, en una sola linea.
\newcommand{\milcestacion}[2]{%
\begin{tcolorbox}[enhanced,colback=#1,colframe=#1,arc=2mm,boxrule=0pt,
  left=5mm,right=5mm,top=2.6mm,bottom=2.6mm,before skip=11pt,after skip=7pt]
{\sffamily\bfseries\fontsize{15}{18}\selectfont\color{white}#2}
\end{tcolorbox}}

% Tarjeta de paso numerada: sustituye las tablas «Pilar / Aplicacion», que
% eran formato de consulta docente, no de estudiante. El numero va en un
% circulo sobre el borde para que la secuencia se lea de un vistazo.
\newcommand{\milcpaso}[3]{%
\begin{tcolorbox}[enhanced,colback=white,colframe=milcLinea,arc=1.5mm,boxrule=.9pt,
  left=14mm,right=4mm,top=3mm,bottom=3mm,before skip=4pt,after skip=4pt,
  fontupper=\RaggedRight,
  overlay={\node[anchor=north west,circle,fill=#2,text=white,inner sep=0pt,
    minimum size=7.4mm,font=\sffamily\bfseries\small]
    at ([xshift=3.6mm,yshift=-3.2mm]frame.north west) {#1};}]
#3
\end{tcolorbox}}

% Casilla marcable de tamano util con lapiz (la anterior era \fbox{\phantom{x}},
% de alto variable segun el texto que la seguia).
\newcommand{\milcmarca}{\framebox[3.8mm]{\rule{0pt}{3.4mm}}}

% Fila marcable de lista suelta (checks de la estacion 1).
\newcommand{\milccheckrow}[1]{%
\par\vspace{1.8mm}\noindent
\makebox[6.4mm][l]{\raisebox{-.55mm}{\textcolor{milcNegro}{\milcmarca}}}%
\parbox[t]{\dimexpr\linewidth-6.4mm\relax}{\RaggedRight #1}\par}

% Celda marcable dentro de la rubrica (tabla de autoevaluacion). Misma
% sangria francesa, pero pensada para vivir dentro de una columna X.
\newcommand{\milcopcion}[1]{%
\makebox[5.2mm][l]{\raisebox{-.55mm}{\textcolor{milcNegro}{\milcmarca}}}%
\parbox[t]{\dimexpr\linewidth-5.2mm\relax}{\RaggedRight #1}}

% ── Recursos visuales de la guía ──────────────────────────────────────
% Declarados en el bloque `recursos:` del YAML y emitidos por el builder.
% \guiaFigura   → imagen rasterizada o PDF (foto, carta celeste, montaje).
% \guiaDiagrama → fuente TikZ/circuitikz (.tex) incrustada como vector:
%                 es la vía para circuitos y esquemáticos editables.
% #1 = ruta relativa al .tex · #2 = pie (puede ir vacío)
\newcommand{\guiaFigura}[2]{%
\begin{center}
\includegraphics[width=0.88\linewidth,height=0.38\textheight,keepaspectratio]{#1}\\[1.5mm]
{\footnotesize\itshape\textcolor{milcNegro!75}{#2}}
\end{center}
}

\newcommand{\guiaDiagrama}[2]{%
\begin{center}
\input{#1}\\[1.5mm]
{\footnotesize\itshape\textcolor{milcNegro!75}{#2}}
\end{center}
}

\begin{document}
\thispagestyle{empty}

% ── Portada ───────────────────────────────────────────────────────────
% Antes: un rectangulo de color a pagina completa. Se veia bien en pantalla,
% pero en la fotocopiadora del colegio se comia el toner y salia gris sucio.
% Ahora la banda ocupa el tercio superior y el resto va en blanco.
\begin{tikzpicture}[remember picture,overlay]
  \path[fill=milcGradePrimary]
    (current page.north west) rectangle ([yshift=-10.6cm]current page.north east);

  \node[anchor=north west,text=milcGradeText] at ([xshift=2.0cm,yshift=-1.55cm]current page.north west)
    {\sffamily\bfseries\fontsize{12}{14}\selectfont MOMENTO};
  \node[anchor=north west,text=milcGradeText,opacity=.85] at ([xshift=2.0cm,yshift=-2.35cm]current page.north west)
    {\sffamily\bfseries\fontsize{8.5}{10}\selectfont TERRITORIO INTERIOR · SOCIOEMOCIONAL};
  \node[anchor=north east,text=milcGradeText,opacity=.85,align=right] at ([xshift=-2.0cm,yshift=-1.55cm]current page.north east)
    {\sffamily\bfseries\fontsize{9}{12}\selectfont I.E. Sor María Juliana\\Cartago, Valle del Cauca};

  \node[anchor=west,text=milcGradeText] (gradonum) at ([xshift=1.85cm,yshift=-7.1cm]current page.north west)
    {\sffamily\bfseries\fontsize{112}{112}\selectfont 9};
  % El circulito de grado (°) solo aplica a las guias de grado («11°»); en
  % Bebras y Semillero el numero grande es un momento/modulo, no un grado.
  % Se posiciona relativo al nodo `gradonum`, no en coordenadas absolutas.
  

  \node[anchor=west,text=milcGradeText,text width=11.4cm,align=left] at ([xshift=1.5cm,yshift=0.5cm]gradonum.east)
    {
     {\sffamily\bfseries\fontsize{9}{11}\selectfont\MakeUppercase{Territorio interior · Socioemocional}}\\[3mm]
     {\sffamily\bfseries\fontsize{21}{25}\selectfont\setlength{\baselineskip}{27pt}Liderar\\[4pt]sin mandar\\[4pt]la autoridad que se recibe}\\[3.5mm]
     {\sffamily\bfseries\fontsize{15}{18}\selectfont Grado 8° · Momento 9}
    };
\end{tikzpicture}

\vspace*{9.4cm}
\vfill

{\sffamily\bfseries\fontsize{8}{10}\selectfont\textcolor{milcGradeDark}{LO QUE TE LLEVAS HOY}}

\begin{tcolorbox}[enhanced,colback=white,colframe=milcNegro,arc=2mm,boxrule=1.6pt,
  left=5mm,right=5mm,top=4mm,bottom=4mm,before skip=3pt,after skip=12pt,
  fontupper=\RaggedRight\fontsize{11}{15.5}\selectfont]
Tu inventario de liderazgo: los liderazgos reales de tu curso clasificados en prestigio o dominancia con la evidencia de cada uno, las tres preguntas de la prueba, y una tarea concreta que vas a coordinar —con el reparto y la fecha de entrega del bastón—.
\end{tcolorbox}

\begin{tcolorbox}[enhanced,colback=milcGris,colframe=milcGris,arc=2mm,boxrule=0pt,
  left=5mm,right=5mm,top=3.5mm,bottom=3.5mm,after skip=12pt]
{\sffamily\bfseries\fontsize{8}{10}\selectfont\textcolor{milcNegro!55}{ESTA GUÍA ES DE}}\\[3mm]
\begin{tabularx}{\linewidth}{X p{3.2cm} p{3.2cm}}
\linead{\linewidth} & \linead{3.0cm} & \linead{3.0cm}\\[-1mm]
{\fontsize{8.5}{10}\selectfont\textcolor{milcNegro!55}{Nombre completo}} &
{\fontsize{8.5}{10}\selectfont\textcolor{milcNegro!55}{Grupo}} &
{\fontsize{8.5}{10}\selectfont\textcolor{milcNegro!55}{Fecha}}\\
\end{tabularx}
\end{tcolorbox}

\vfill

\noindent\textcolor{milcLinea}{\rule{\linewidth}{1pt}}\\[2mm]
{\sffamily\bfseries\fontsize{9.5}{12}\selectfont Autor: Dr. Álvaro Cárdenas Orozco}\\[1mm]
{\sffamily\fontsize{8.5}{11}\selectfont\textcolor{milcNegro!55}{Metodología MILC · Apertura ancestral · Autoridad y cuidado · Triángulo Dussel-Estoicismo-Floridi}}

\newpage

\section*{Apertura: Liderar sin mandar: la autoridad que se recibe y la que se impone}

\begin{softbox}{milcGradeDark}{milcGradeSoft}{Saber ancestral}
En los pueblos indígenas de Colombia, la autoridad tiene \textbf{un objeto que se ve}: el \textbf{bastón de mando}. Lo llevan las autoridades del \textbf{cabildo}, y no es un adorno ---es la señal de que esa persona está ejerciendo un cargo---. \textbf{Y lo importante no es el bastón: es cómo se mueve.} \textbf{Primero: no se compra ni se toma. Se recibe.} \emph{Es la comunidad la que lo entrega, y con él entrega el encargo.} \textbf{Segundo: es por un periodo.} Los cabildos \textbf{se posesionan} y se renuevan; \emph{cuando entra un cabildo nuevo, los bastones se refrescan} ---la autoridad se limpia y vuelve a empezar---. \textbf{Tercero, y es lo que más cuesta entender a los trece años: el bastón se devuelve.} \emph{Nadie es autoridad para siempre, y dejar el cargo no es un fracaso: es parte del cargo.} \textbf{Y hay una cuarta cosa, la decisiva:} \emph{lo que el bastón entrega no es el derecho a mandar, sino \textbf{la obligación de responder por otros}} ---por el territorio, por la comunidad, por lo que pase durante ese periodo---. \textbf{Autoridad, ahí, es una carga que se presta.}
\end{softbox}

\begin{softbox}{milcTurquesa}{milcGris}{Saber contemporáneo}
¿Cómo llega alguien a tener influencia sobre un grupo? \textbf{Hay dos caminos, y llevan al mismo puesto por rutas opuestas.} \textbf{Joseph Henrich} y \textbf{Francisco Gil-White} los describieron así: la \textbf{dominancia} es el uso de \textbf{la coerción y la amenaza para infundir miedo}; el \textbf{prestigio}, en cambio, es una influencia \textbf{libremente concedida}, que además trae beneficios para los demás (Henrich y Gil-White, 2001). \textbf{Y su explicación del prestigio es lo mejor de todo:} la gente le concede deferencia a alguien \textbf{porque quiere aprender de él}. \emph{Somos aprendices ---copiamos a los que parecen saber más--- y por eso nos acercamos a esas personas, las escuchamos y les damos un lugar.} \textbf{Fíjate en la palabra clave: \emph{libremente concedida}.} \emph{El prestigio no se toma: te lo dan, y te lo pueden quitar.} \textbf{Ahora aplica los dos a tu salón.} El que manda por \textbf{dominancia} consigue que le \textbf{obedezcan cuando está mirando}; el que tiene \textbf{prestigio} consigue que lo \textbf{busquen cuando no está}. \emph{Desde afuera los dos parecen «los líderes». Por dentro no se parecen en nada ---y se distinguen con una sola pregunta: qué pasa cuando esa persona no está---.}
\end{softbox}

\begin{softbox}{milcMostaza}{milcGris}{Pregunta-puente}
\textit{El bastón \textbf{lo entrega la comunidad, es por un periodo y se devuelve}. \textbf{¿A quién le harían caso en tu curso si mañana no estuviera el que manda hoy}\ldots \textbf{y por qué a esa persona}?}
\end{softbox}

\begin{softbox}{milcVerde}{milcGris}{Saber y saber hacer}
Al terminar podrás: (1) \textbf{distinguir} la dominancia del prestigio con evidencia observable, y no con impresiones; (2) \textbf{explicar} por qué el miedo consigue obediencia pero no seguimiento; (3) \textbf{entender} la autoridad como algo que se recibe, se ejerce por un periodo y se devuelve; (4) \textbf{coordinar} una tarea real repartiendo el trabajo en vez de acumularlo, y cerrarla entregando el puesto.
\end{softbox}



\small
\begin{tabularx}{\textwidth}{>{\bfseries}p{3.0cm}Y}
\rowcolor{milcGradePrimary!12}Autor & Dr. Álvaro Cárdenas Orozco\\
DBA & Comprende los fenómenos que afectan a su territorio, regula sus emociones ante ellos y acompaña a otros con empatía (transversal 6°--11°, articulado con la política «Escuela, territorio de vida» del MEN, Resolución 006519 de 2025, y la Ley 1620 de 2013)\\
Referentes & Primera ayuda psicológica (OMS, 2011/2012) · Directrices IASC (2007) · USGS y Servicio Geológico Colombiano · Pueblo nasa y la minga (CRIC) · Dussel · Estoicismo · Floridi\\
Producto final & Tu inventario de liderazgo: los liderazgos reales de tu curso clasificados en prestigio o dominancia con la evidencia de cada uno, las tres preguntas de la prueba, y una tarea concreta que vas a coordinar —con el reparto y la fecha de entrega del bastón—.\\
\end{tabularx}
\normalsize

\puente{Ya viste cómo se fabrican las famas, las varas y las cercas. \textbf{Hoy: cómo se fabrica la autoridad} ---y por qué hay dos maneras muy distintas de conseguirla---.}

\milcestacion{milcGradeDark}{Ruta MILC: así vamos a aprender}
\begin{tabularx}{\textwidth}{>{\centering\arraybackslash}X>{\centering\arraybackslash}X>{\centering\arraybackslash}X>{\centering\arraybackslash}X}
\phasecard{1}{milcVerde}{milcGris}{Escucha\\[-1mm]\small Observo, escucho y pregunto.} &
\phasecard{2}{milcTurquesa}{milcGris}{Sistematización\\[-1mm]\small Distingo y sirvo.} &
\phasecard{3}{milcMagenta}{milcGris}{Praxis\\[-1mm]\small Construyo y aplico.} &
\phasecard{4}{milcVino}{milcGris}{Evaluación\\[-1mm]\small Reflexión liberadora.}
\end{tabularx}


\puente{Primero, el inventario real: quién manda en tu curso, quién arrastra y quién es buscado cuando hay un problema. \emph{No siempre son la misma persona.}}

\milcestacion{milcVerde}{Estación 1 de 3 · Escucha}

\begin{softbox}{milcVerde}{milcGris}{Escena para escuchar}
\textbf{Actividad 1 · IDENTIFICA --- Quién manda y quién es buscado.} \textbf{Sin nombres: usa iniciales.} \textbf{Primera parte.} Responde tres preguntas distintas sobre tu curso: \emph{¿Quién decide lo que hace el grupo? ¿A quién le hacen caso? ¿A quién buscan cuando hay un problema de verdad ---algo perdido, alguien llorando, un lío con un profesor---?} \textbf{Fíjate si son la misma persona o no.} \textbf{Segunda parte.} Para cada una escribe \textbf{cómo consiguió ese lugar}: ¿porque sabe algo? ¿porque ayuda? ¿porque da miedo quedar mal con ella? ¿porque tiene con qué? \textbf{Tercera parte, la reveladora.} \emph{Cuando esa persona no está, ¿qué pasa?} \textbf{Cuarta.} ¿En qué eres tú \textbf{el que otros buscan}? \emph{Aunque sea en algo pequeño ---prestar algo, explicar una materia, saber una fecha---: casi todo el mundo es referencia de alguien y casi nadie lo sabe.}
\end{softbox}

{\sffamily\bfseries\fontsize{8}{10}\selectfont\textcolor{milcNegro!55}{MARCA CUANDO LO HAYAS HECHO}}
\milccheckrow{Respondiste las tres preguntas y notaste si son o no la misma persona.}
\milccheckrow{Escribiste \textbf{cómo consiguió} cada una su lugar.}
\milccheckrow{Anotaste qué pasa \textbf{cuando esa persona no está}.}

\begin{infoband}{milcVerde}


\textbf{En tu cuaderno (Actividad 1 · IDENTIFICA).} Título: «Actividad 1. Quién manda, quién es buscado». \textbf{Formato:} las tres preguntas con iniciales + cómo consiguió cada uno su lugar + qué pasa cuando no está + en qué eres tú referencia. \textbf{Extensión:} media página. \textbf{Sabes que terminaste cuando} puedes decir si \textbf{el que manda y el que es buscado son la misma persona}. Esta página es tuya: \textbf{nadie más tiene que leerla si tú no quieres}.
\end{infoband}

\puente{Ya lo tienes. Ahora los dos caminos ---dominancia y prestigio---, y qué le pasa a cada uno cuando el líder no está.}

\milcestacion{milcTurquesa}{Fase 2 · Sistematización: entender la autoridad}

\begin{softbox}{milcTurquesa}{milcGris}{Cuatro claves sobre liderar}
Cuatro claves para dejar de confundir al que manda con el que lidera.
\end{softbox}

\milcpaso{1}{milcTurquesa}{\textbf{Dos caminos al mismo puesto.} \textbf{Dominancia:} coerción y amenaza, \textbf{para infundir miedo}. \textbf{Prestigio:} influencia \textbf{libremente concedida}, que trae beneficios a los demás; se la damos a quien creemos que sabe algo, \textbf{porque queremos aprender de él} (Henrich y Gil-White, 2001). \emph{Desde afuera los dos se ven igual ---los dos tienen gente alrededor---.} \textbf{La prueba que los separa es una sola:} \emph{¿qué pasa cuando esa persona no está?} \textbf{Al que tiene prestigio lo siguen buscando; al que domina, lo dejan de obedecer en el mismo recreo.} \emph{Porque el prestigio se lo dieron ---y eso no se va cuando él se va---, y la dominancia se sostiene únicamente mientras esté mirando.}}
\milcpaso{2}{milcTurquesa}{\textbf{El bastón se entrega, no se toma.} En el cabildo, la autoridad \textbf{la da la comunidad}, dura \textbf{un periodo} y \textbf{se devuelve}. \emph{Es exactamente lo que Henrich y Gil-White dicen del prestigio con otras palabras:} \textbf{es \emph{libremente concedido}}, y lo que se concede se puede retirar. \emph{Y hay una consecuencia práctica para cualquiera que quiera liderar algo:} \textbf{no se puede \emph{tomar} el liderazgo de un grupo}. \emph{Lo único que está en tu poder es hacer cosas por las que el grupo quiera dártelo} ---y después, sostenerlo---. \textbf{Y otra, que casi nadie enseña:} \emph{dejar el cargo cuando toca es parte de tenerlo}. \textbf{En el cabildo eso no es una derrota: los bastones se refrescan y entra el siguiente.}}
\milcpaso{3}{milcTurquesa}{\textbf{El miedo consigue obediencia, no seguimiento.} \emph{La dominancia funciona ---esa es la parte incómoda---:} en un salón, el que da miedo consigue que le presten, que le copien, que le guarden el puesto. \textbf{Pero fíjate en el precio, que se paga entero y en silencio:} \emph{nadie le cuenta nada verdadero}. \textbf{Al que domina no le avisan cuando se está equivocando}, no le dicen que está mal, no le piden ayuda ---\emph{y por eso el que manda por miedo casi siempre es el último en enterarse de todo lo que pasa en su propio grupo}---. \textbf{Y hay algo más:} \emph{la dominancia se tiene que reafirmar}. \textbf{El prestigio, no.}}
\milcpaso{4}{milcTurquesa}{\textbf{Liderar es repartir, no acumular.} El error más común de un buen estudiante que «coordina» un trabajo es \textbf{hacerlo casi todo}, porque le sale más rápido y le queda mejor. \emph{Eso no es coordinar: es trabajar solo con público.} \textbf{Coordinar es otra cosa y es más difícil:} \emph{repartir según lo que cada uno puede} ---exactamente como en la minga, donde a cada quien se le da una tarea a su medida---, \textbf{revisar sin humillar y hacer que el trabajo del otro se note}. \emph{Y esta es la prueba más honesta de todas:} \textbf{si el grupo puede terminar sin ti, coordinaste bien}.}

\begin{drawbox}{milcTurquesa}{La prueba de las tres preguntas}
\begin{pasoslista}
\item \textbf{¿Qué pasa cuando esa persona no está?} \emph{Si todo se cae, era dominancia. Si sigue, era prestigio.}
\item \textbf{¿La gente le cuenta cosas verdaderas?} \emph{Al que da miedo no le avisan cuando se está equivocando.}
\item \textbf{¿Reparte o acumula?} \emph{Mira si los demás tienen algo que hacer, o solo obedecen.}
\item \textbf{Y una cuarta, para ti:} \emph{¿estoy tomando un puesto, o haciendo algo por lo que me lo quieran dar?}
\end{pasoslista}
\end{drawbox}

\begin{softbox}{milcMostaza}{milcGris}{Errores comunes y su corrección}
\textbf{Confundir mandar con liderar:} el que manda tiene obediencia; el que lidera tiene gente que lo busca. \textbf{Hacerlo todo uno:} eso es trabajar con público. \textbf{Creer que el liderazgo se toma:} se concede, y se puede retirar. \textbf{Revisar humillando:} corrige el trabajo y pierde a la persona. \textbf{No soltar el puesto:} el bastón se devuelve; quedarse pegado lo convierte en dominancia. \textbf{Repartir parejo sin mirar a quién:} en la minga se reparte según lo que cada uno puede. \textbf{Quedarse con el crédito:} el que se lleva el reconocimiento del grupo lo pierde a la segunda vez.
\end{softbox}

\begin{infoband}{milcTurquesa}


\textbf{En tu cuaderno (Actividad 2 · EXPLICA).} Título: «Actividad 2. Dominancia y prestigio». \textbf{Formato:} dos columnas con las características de cada camino, y clasifica ---sin nombres--- los tres liderazgos de la Actividad 1, con \textbf{la evidencia} en que te basas. \textbf{Extensión:} media página. \textbf{Sabes que terminaste cuando} puedes explicar por qué \textbf{al que domina casi nadie le cuenta la verdad}.
\end{infoband}

\puente{Ya sabes distinguirlos. Es hora de ejercerlo: coordinar algo de verdad, que es donde se aprende esto y no leyendo.}

\milcestacion{milcMagenta}{Fase 3 · Praxis: coordinar algo}

\begin{softbox}{milcMagenta}{milcGris}{Cómo se coordina sin mandar}
Esto no se aprende leyéndolo. Se aprende coordinando algo pequeño de verdad, y descubriendo lo difícil que es no hacerlo todo uno mismo.
\end{softbox}

\milcpaso{1}{milcMagenta}{\textbf{El inventario, clasificado (8 min).} Pasa los tres liderazgos de tu curso a las dos columnas, \textbf{con evidencia observable} ---«se le nota que\ldots» no es evidencia; «cuando faltó, el grupo siguió» sí lo es---. \emph{Y ponte tú también en la tabla, aunque sea en algo mínimo.}}
\milcpaso{2}{milcMagenta}{\textbf{La prueba de las tres preguntas (7 min).} Aplícala al liderazgo más fuerte de tu curso. \textbf{Presta atención a la segunda} ---\emph{¿le cuentan cosas verdaderas?}---: \emph{es la que más rápido revela de qué tipo es, y la que casi nadie mira}.}
\milcpaso{3}{milcMagenta}{\textbf{La tarea que voy a coordinar (12 min).} Escoge \textbf{algo real y pequeño} ---un trabajo de clase, el aseo de la semana, un partido, un montaje para la izada, una parte del proyecto del momento 10--- y escribe: \textbf{qué es, quiénes participan, qué le toca a cada uno y por qué a esa persona le toca eso}. \emph{Regla de la minga:} \textbf{a cada quien según lo que puede hacer}, y \textbf{nadie sin tarea}. \emph{Y una condición que va a doler:} \textbf{tú no puedes tener la tarea más grande.}}
\milcpaso{4}{milcMagenta}{\textbf{La entrega del bastón (5 min y al final).} Escribe desde ahora \textbf{cuándo termina tu coordinación} y \textbf{cómo la vas a cerrar}: qué se entrega, a quién, y ---esto es lo importante--- \textbf{cómo vas a hacer que se note lo que hizo cada uno}. \emph{Al final, haz dos cosas:} \textbf{decir en voz alta qué hizo cada quien} y \textbf{devolver el puesto} ---que el siguiente lo coordine otro---. \emph{Esa última parte es la que convierte esto en un bastón y no en un trono.}}

\begin{drawbox}{milcMagenta}{Antes de cerrar tu inventario}
\begin{checklista}
\item Clasificaste los liderazgos con \textbf{evidencia observable}, no con impresiones.
\item Te incluiste a ti en la tabla.
\item Aplicaste las tres preguntas, incluida la de si le cuentan cosas verdaderas.
\item Tu tarea repartida deja a \textbf{nadie sin trabajo} y a ti \textbf{sin la parte más grande}.
\item Escribiste cuándo termina tu coordinación y cómo la cierras.
\item Planeaste \textbf{cómo se va a notar lo que hizo cada uno} ---y que el próximo lo coordine otro---.
\end{checklista}
\end{drawbox}

\begin{softbox}{milcMagenta}{milcGris}{Plantilla de trabajo}
\textbf{Las tres preguntas.} \emph{«¿Qué pasa cuando no está? · ¿Le cuentan cosas verdaderas? · ¿Reparte o acumula?»} \\[1mm]
\textbf{Mi reparto.} \emph{«\_\_\_\_ hace \_\_\_\_ porque \_\_\_\_.»} ---uno por cada participante, nadie sin tarea---. \\[1mm]
\textbf{La entrega.} \emph{«Esto termina el \_\_\_\_. Lo hecho por cada uno fue \_\_\_\_. El próximo lo coordina \_\_\_\_.»}
\end{softbox}

\begin{infoband}{milcMagenta}


\textbf{En tu cuaderno (Actividad 3 · CREA).} Título: «Actividad 3. Lo que voy a coordinar». \textbf{Formato:} el inventario clasificado con evidencia + las tres preguntas + la tarea con el reparto justificado + la fecha y la forma de entrega del bastón. \textbf{Extensión:} una página. \textbf{Sabes que terminaste cuando} en tu reparto \textbf{tú no tienes la parte más grande}.
\end{infoband}

\puente{El producto es una coordinación con fecha de entrega. \emph{La parte más difícil no es organizar: es soltar el puesto al final.}}

\milcestacion{milcMostaza}{Producto de la sesión}
\textbf{Inventario de liderazgo: prestigio y dominancia con evidencia, y una tarea coordinada con entrega del bastón}.
\begin{softbox}{milcMostaza}{milcGris}{Criterios de calidad}
(1) La clasificación se apoya en evidencia observable, no en impresiones sobre las personas. (2) Se aplica la prueba de las tres preguntas, incluida la de si al líder le cuentan la verdad. (3) El reparto asigna tarea a todos, según lo que cada uno puede hacer, y no concentra el trabajo en quien coordina. (4) La coordinación tiene fecha de cierre y una forma explícita de hacer visible el aporte de cada uno. (5) Se devuelve el puesto al terminar: el siguiente turno lo coordina otra persona.
\end{softbox}

\puente{Antes de cerrar, revisa el error más común: si terminaste haciendo casi todo tú, no coordinaste ---trabajaste---.}

\milcestacion{milcVino}{Evaluación liberadora · cinco dimensiones}

\begin{softbox}{milcVino}{milcGris}{Sentido de la evaluación}
Evaluar no es solo revisar una nota. Es reconocer qué aprendí, qué cambió en mí, qué emociones manejé, cómo me relaciono con la ciudad y la comunidad, y qué le dejaré a quien viene detrás.
\end{softbox}

\textbf{Mi reflexión liberadora en cinco dimensiones.} Completa cada frase iniciada:

\small
\begin{tabularx}{\textwidth}{>{\bfseries\RaggedRight\arraybackslash}p{3.4cm}Y>{\RaggedRight\arraybackslash}p{4.6cm}}
\rowcolor{milcVino}\color{white}Dimensión & \color{white}\bfseries Mi respuesta (frame starter) & \color{white}\bfseries Algo que descubrí\\
1. Personal & Lo que más cambió en mí fue \linead{6.5cm} & \linead{4.2cm}\\
2. Emocional & La emoción más fuerte fue \linead{2.5cm} y la regulé \linead{3cm} & \linead{4.2cm}\\
3. Ciudadana & En democracia esto importa porque \linead{6cm} & \linead{4.2cm}\\
4. Local (Cartago) & En mi barrio sería útil para \linead{6.5cm} & \linead{4.2cm}\\
5. Intergeneracional & A mi abuela/abuelo le diría \linead{6.5cm} & \linead{4.2cm}\\
\end{tabularx}
\normalsize

\begin{infoband}{milcVino}
\textbf{Para la próxima sustentación.} Vuelve a esta tabla en seis meses. Verás que las respuestas cambian — no porque lo aprendido cambie, sino porque tú cambias. La evaluación liberadora es una conversación contigo a través del tiempo.
\end{infoband}


\puente{Cerramos con tres voces: la comunidad que se funda escuchando, la coherencia entre palabras y acciones, y el cuidado de lo compartido.}

\milcestacion{milcGradeDark}{Triángulo de pensamiento · cierre filosófico}

\begin{softbox}{milcVino}{milcGris}{Dussel · lente del nosotros}
\textit{«Aquel que es capaz de escuchar silenciosamente al Otro como otro es el que puede constituir una comunidad y no una mera sociedad totalizada.»}

{\footnotesize\itshape\color{milcNegro!70}--- Enrique Dussel · Filosofía de la liberación (1977)}

Dussel pone la escucha en el origen de la comunidad, \textbf{y eso es una definición de liderazgo mucho mejor que cualquier lista de cualidades}. \emph{El que manda por dominancia no puede escuchar ---y no porque no quiera: es que nadie le dice nada verdadero---.} \textbf{Se queda encerrado en su propia versión del grupo}, y por eso los grupos que dirige funcionan bien mientras nada se salga del libreto. \emph{Fíjate en la conexión con el primer momento del grado 7:} \textbf{escuchar sin arreglar no era solo una habilidad para consolar amigos ---es la herramienta central del que coordina algo---}. \emph{Nadie puede repartir bien un trabajo si no sabe qué puede hacer cada quien, y eso solo se averigua preguntando y aguantando la respuesta.}

\textbf{Cuando organizo algo, ¿pregunto qué puede hacer cada uno, o reparto según lo que yo supongo?}
\end{softbox}
\linead{\linewidth}\\[1mm]
\linead{\linewidth}

\begin{softbox}{milcMostaza}{milcGris}{Estoicismo · Séneca · Cartas a Lucilio, 20 (c. 64 d.C.) · cuidado interior}
\textit{«La filosofía enseña a obrar, no a hablar; quiere que cada uno viva a la manera que prescribe, que estén en armonía nuestras palabras con nuestras acciones.»}

{\footnotesize\itshape\color{milcNegro!70}--- Séneca · Cartas a Lucilio, 20 (c. 64 d.C.)}

\textbf{«Que estén en armonía nuestras palabras con nuestras acciones».} \emph{Esa frase es, exactamente, de dónde sale el prestigio.} \textbf{Henrich y Gil-White dicen que la gente concede deferencia a quien cree que sabe algo ---y «saber» aquí no se demuestra hablando---:} \emph{se demuestra haciendo, donde los demás lo vean}. \textbf{Por eso el liderazgo de un curso casi nunca lo gana el que más habla:} \emph{lo gana el que cumplió lo que dijo, dos o tres veces seguidas, incluso cuando le costaba}. \textbf{Y por eso también se pierde tan rápido:} \emph{una sola vez que digas que vas a hacer algo y no lo hagas, el grupo lo registra ---y como viste en el momento 4, lo malo pesa más---.}

\textbf{¿Cuál fue la última cosa que dije que iba a hacer y no hice? ¿Quién se dio cuenta?}
\end{softbox}
\linead{\linewidth}\\[1mm]
\linead{\linewidth}

\begin{softbox}{milcTurquesa}{milcGris}{Floridi · lente de la infoesfera}
\textit{«El florecimiento de las entidades informacionales y de toda la infoesfera debe promoverse preservando, cultivando y enriqueciendo sus propiedades.»}

{\footnotesize\itshape\color{milcNegro!70}--- Luciano Floridi · La ética de la información (2013; trad.)}

\textbf{Cultivar} es la palabra exacta para lo que hace alguien que coordina bien: \emph{no ocupa el espacio ---lo prepara para que otros hagan---}. \textbf{Y hay una versión muy concreta en cualquier grupo de trabajo hoy:} \emph{quien administra el grupo de chat tiene un poder que casi nadie nombra} ---decide quién está, qué se responde, qué se deja pasar---. \textbf{Ese también es un bastón, y también se puede usar de las dos maneras:} \emph{para que todos se enteren, o para que algunos no}. \textbf{La prueba es la misma de siempre:} \emph{si el grupo funciona igual de bien cuando tú no estás mirando, lo estabas cultivando; si se apaga, lo estabas ocupando.}

\textbf{En los grupos donde tengo algún papel, ¿las cosas siguen cuando yo no estoy?}
\end{softbox}
\linead{\linewidth}\\[1mm]
\linead{\linewidth}

\begin{infoband}{milcNegro}
\textbf{Nota para el docente.} Las tres son citas textuales y llevan su referencia debajo. Puedes remitir a la obra en clase.
\end{infoband}

\milcestacion{milcVino}{Mi autoevaluación · marca una casilla por fila}

{\small
\setlength{\extrarowheight}{2.2pt}
\begin{tabularx}{\textwidth}{>{\RaggedRight\arraybackslash\bfseries}p{3.0cm}YYY}
\rowcolor{milcVino}\color{white}\bfseries Criterio & \color{white}\bfseries Lo logré & \color{white}\bfseries En proceso & \color{white}\bfseries Necesito apoyo\\[.6mm]
Comprensión del tema & \milcopcion{Explico las ideas con mis palabras.} & \milcopcion{Reconozco las ideas, falta precisión.} & \milcopcion{Necesito repasar lo básico.}\\[.8mm]
\rowcolor{milcGris}Calidad de la coordinación & \milcopcion{Coordiné una tarea repartiendo el trabajo y devolví el puesto al terminar.} & \milcopcion{Organicé algo, pero terminé haciéndolo casi todo yo.} & \milcopcion{Sigo creyendo que liderar es que a uno le hagan caso.}\\[.8mm]
Aplicación y producto & \milcopcion{Mi producto cumple los criterios.} & \milcopcion{Está armado, con ajustes pendientes.} & \milcopcion{Aún está en construcción.}\\[.8mm]
\rowcolor{milcGris}Reflexión 5 dimensiones & \milcopcion{Respondí cada una con honestidad.} & \milcopcion{Respondí algunas con detalle.} & \milcopcion{Mis respuestas son muy generales.}\\[.8mm]
Triángulo de pensamiento & \milcopcion{Conecté las 3 voces con mi trabajo.} & \milcopcion{Conecté al menos una.} & \milcopcion{No las relaciono con mi proceso.}\\[.8mm]
\end{tabularx}}

\vspace{3mm}
\textbf{Hoy entendí que:}\\[1mm]
\linead{\linewidth}\\[1mm]
\linead{\linewidth}

\textbf{Mi compromiso para la próxima sesión es:}\\[1mm]
\linead{\linewidth}\\[1mm]
\linead{\linewidth}

\begin{softbox}{milcMostaza}{milcGris}{Cita de despedida · Marco Aurelio}
\textit{``El obstáculo en el camino se vuelve el camino. Lo que estorba al camino se vuelve camino.''}\\[1mm]
Cada error de hoy, bien observado, se convierte en pieza del camino que pisarás mañana.
\end{softbox}

\small
\begin{tabularx}{\textwidth}{>{\bfseries}p{3.6cm}Y}
\rowcolor{milcGradeDark!12}Nombre completo: & \linead{10cm}\\
Fecha: & \linead{4cm} \quad Curso/grupo: \linead{4cm}\\
Mi firma: & \linead{9cm}\\
\end{tabularx}
\normalsize

\begin{infoband}{milcGradeDark}
\textbf{Para la próxima sesión.} Dos cosas. \textbf{Primero, coordina la tarea} ---y trae dos datos concretos: \textbf{qué hizo cada uno} y \textbf{cuánto hiciste tú}---. \emph{Si al final resulta que hiciste más de la mitad, escríbelo: es el hallazgo más común y el más útil de esta guía.} \textbf{Segundo, pregunta en tu casa quién era el que la gente buscaba:} averigua con alguien mayor \textbf{quién era la persona a la que acudían en su barrio o su vereda cuando había un problema} ---no la de más plata ni la del cargo: la que la gente buscaba---, \textbf{por qué era esa} y qué hacía. \textbf{Trae:} los dos datos de tu coordinación y lo que te contaron. \emph{Vas a encontrar casi siempre lo mismo, y encaja con todo lo que vimos hoy: no era la que más mandaba. Era la que había resuelto cosas y la que sabía escuchar ---y nadie la había nombrado para eso---.}
\end{infoband}

\end{document}
